\documentclass[12pt]{article}
\usepackage[utf8]{inputenc}
\usepackage[T1]{fontenc}
\usepackage{hyperref}
\title{Koraga, Mudu Language Resources for K-12 Curriculum Development}
\author{Compiled by Claude Sonnet 5 (claude-sonnet-5) for Andr\'as Kornai}
\date{}
\begin{document}
\maketitle
\section{Introduction}
Mudu Koraga (ISO 639-3 vmd, Glottocode mudu1241) is a Dravidian language of India. Wikipedia gives roughly 45--50 native speakers (2018), all bilingual in Kannada. It is sometimes treated as a dialect of Koraga.
\section{Primary Educational Resources}
None found.
\section{Linguistic and Reference Materials}
General Koraga article: \url{https://en.wikipedia.org/wiki/Koraga_language}. Endangered Languages entry and its bibliography (citing van Driem's chapter in the \emph{Encyclopedia of the World's Endangered Languages}, 2007): \url{https://endangeredlanguages.com/lang/5651} and \url{https://endangeredlanguages.com/lang/5651/bibliography}. Glottolog: \url{https://glottolog.org/resource/languoid/id/mudu1241.bigmap.html}. A bookseller listing for a Koraga monograph: \url{https://www.exoticindiaart.com/book/details/NAM053/}. No grammar, dictionary or corpus specific to Mudu was found.
\section{Religious and Community Resources}
None found.
\section{Digital and Audio Resources}
None found.
\section{Curriculum Development Context}
Very few speakers remain and documentation is thin; a curriculum would need new documentation work with the community, with Kannada as the contact language.
\end{document}